\section{扩展话题：多任务与分布漂移}
\subsection{多任务 policy gradient}
在 multi-task RL 中，policy gradient 需要在不同任务间共享参数。最常见做法是使用 shared encoder + task-specific head，这允许在不同 reward function 之间快速切换。训练时通常加入 task-id embedding 以区分 contexts。

\begin{knowledgebox}{多任务训练要点}
\begin{itemize}
    \item 同一 batch 中混合多个任务，若 reward scale 不同需标准化；
    \item 使用 task embedding 让 policy aware of active task；
    \item Critic 可以共享 backbone，但多头 value head 用于分别估计 advantage。
\end{itemize}
\end{knowledgebox}

\subsection{分布漂移与 domain adaptation}
在实际部署时，训练数据与线上数据之间的差异会导致 policy drift。常见缓解方式包括：加入 domain randomization、使用 adversarial buffer 采样 hard cases、在 rollout 中自动扩充 replay buffer（cohort of policies）。我们需在 log 中明确记录 `data shift ratio` 与 `reward shift ratio`。

\begin{warningbox}{忽略 domain shift 的后果}
忽略分布漂移会让策略在小改动（如 sensor noise）下性能崩溃，纪要到 compliance review 可能会被判定为 “lacks robustness”。
\end{warningbox}

\subsection{本章小结}
多任务训练与分布漂移是延展策略梯度到实际系统后的关键挑战，需在 policy architecture + logging pipeline 上做足手脚。

